\subsection{¿Cómo afectan los distintos $w(t)$ y la elección de $q$ al desempeño de $S_n$?}

En primer lugar, respecto a la elección de \(w(t)\), se observa que las tres alternativas consideradas producen comportamientos similares. Ya que la potencia aumenta conforme crece el tamaño muestral \(n\), alcanzando valores cercanos a 1 en la mayoría de escenarios. Sin embargo, sí aparecen diferencias relevantes en la rapidez con que esto ocurre.

La ponderación Laplace muestra una convergencia más rápida hacia potencia alta, especialmente frente al escenario Pareto. Esto sugiere que asignar mayor peso relativo a regiones alejadas del origen permite capturar mejor desviaciones de simetría asociadas a colas asimétricas. En contraste, los pesos gaussianos tienden a concentrar mayor masa alrededor de \(t=0\), privilegiando información local de baja frecuencia. Esto produce un comportamiento algo más estable, aunque en algunos escenarios requiere tamaños muestrales ligeramente mayores para alcanzar el mismo nivel de potencia.

Entre los pesos gaussianos, la diferencia entre \(\sigma=0.5\) y \(\sigma=1.0\) es moderada. El caso \(\sigma=0.5\) concentra aún más la integral cerca del origen, mientras que \(\sigma=1.0\) distribuye el peso sobre una región más amplia. Empíricamente no se observan diferencias dramáticas entre ambas, aunque \(w=\mathrm{gauss}_{1.0}\) parece ser más estable.

Por otra parte, el parámetro \(q\) controla la forma de penalización de las discrepancias. Cuando \(q=1\), el estadístico agrega desviaciones de manera aproximadamente lineal, resultando menos sensible a diferencias extremas puntuales. Esto genera un comportamiento más robusto y estable frente a variabilidad muestral. Cuando \(q=2\), las discrepancias grandes reciben una penalización cuadrática, aumentando la sensibilidad del test frente a desviaciones pronunciadas respecto a simetría.

En términos prácticos, esto implica que \(q=2\) puede detectar con mayor facilidad asimetrías marcadas, mientras que \(q=1\) tiende a comportarse mejor cuando la asimetría está distribuida de forma más homogénea a lo largo de la distribución. En las simulaciones $q=1$ resulta sistemáticamente superior a $q=2$ con mediana y media afeitada (por ejemplo, con media afeitada y Pareto a $n=40$ la potencia pasa de $38.5\%$ con $q=2$ a $50\%$ con $q=1$). No obstante, este aumento de sensibilidad también puede amplificar el ruido numérico proveniente de la estimación de la función característica empírica, especialmente para valores grandes de \(t\).

Respecto a la elección del estimador del centro $\hat\theta$, los resultados son matizados (véase el Apéndice). Bajo Cauchy, la mediana exhibe el menor RMSE como estimador puntual, mientras que bajo Uniforme el orden se invierte: el argmin tiene el menor RMSE, y la mediana el mayor. En cuanto a la potencia de $S_n$, el argmin es de hecho \emph{el estimador más potente} cuando es compatible con $q$: para Gamma con $q=2$ y $w_1$, alcanza $47.5\%$ en $n=40$ frente al $32\%$ de la mediana; ventaja que se mantiene en todas las distribuciones alternativas. La explicación es que el argmin minimiza $S_n$ sobre la muestra observada, dejando $S_n^{\text{obs}}$ pequeño pero, simultáneamente, las remuestras bootstrap (provenientes de $sF_n(\hat\theta_{\min})$) producen $S_n^*$ comparablemente pequeño, de manera que la calibración relativa es correcta y el test detecta la asimetría con mayor sensibilidad. La limitación práctica del argmin es su exclusión bajo $q=1$ —ya que el integrando $|\cdot|^{1/2}$ no es diferenciable en sus ceros y la búsqueda en grilla resulta prohibitiva en el bootstrap— por lo que en ese caso solo se evalúan mediana y media afeitada.

De forma conjunta, los resultados sugieren que la elección de \(w(t)\) afecta principalmente la sensibilidad del test frente a distintos tipos de asimetría —especialmente en colas— mientras que \(q\) modula la severidad con que dichas discrepancias son penalizadas. Empíricamente, las configuraciones con peso Laplace o gaussiano de escala intermedia, combinadas con el argmin (cuando $q=2$) o con la media afeitada (cuando $q=1$), presentan el mejor compromiso entre potencia y estabilidad numérica.
